% Actual class: 03
% Slide decks: CZ3005 Module 3_Constraint satisfaction and adversarial search; CZ3005 Module 4_Markov Decision Process
% Exact scope: Module 3 pp. 26--54; Module 4 pp. 1--24, aligned with the Week 4 schedule
% NTULearn supplement: None
% Human verified: Concepts discussed and checks completed

\section{第 03 节课：Adversarial Search 与 MDP 1}
\label{lecture:SC3000-L03}

\lecturemeta
  {SC3000-L03}
  {2026-08-31}
  {CZ3005 Module 3；CZ3005 Module 4}
  {Module 3 pp.~26--54；Module 4 pp.~1--24（对应 Week 4 schedule）}
  {无}
  {知识内容已讨论并通过检查题}

\subsection{本讲主线}

Adversarial search（对抗搜索）面对会主动选择不利动作的 opponent（对手），用 Minimax（极大极小算法）在 MAX 与 MIN 之间回传 utility（效用）。Markov Decision Process（马尔可夫决策过程，MDP）面对 stochastic transition（随机转移），用 probability-weighted expectation（概率加权期望）评价动作。两者都研究 sequential decision making（序贯决策），但不确定性的来源和备份运算不同。

\lecturetopic{SC3000-L03-T01}{Game Formulation 与 Minimax}

本讲聚焦 two-player、turn-taking、deterministic、fully observable、zero-sum games。一个 game problem 由 initial state、current player、legal actions、result function、terminal test 与 utility function 构成；utility 统一从 MAX 视角定义，例如胜、和、负分别记为 $+1,0,-1$。

Minimax 假设双方都 optimal（最优）：MAX 选最大值，MIN 选最小值。
\[
  V(s)=
  \begin{cases}
    U(s), & s\text{ 为 terminal state},\\
    \displaystyle\max_{a\in A(s)}V(\operatorname{Result}(s,a)),
      & s\text{ 为 MAX node},\\
    \displaystyle\min_{a\in A(s)}V(\operatorname{Result}(s,a)),
      & s\text{ 为 MIN node}.
  \end{cases}
\]
Minimax value 是面对 optimal opponent 时能够保证的结果；若 opponent 犯错，实际收益可以更高。

\begin{figure}[htbp]
  \centering
  \resizebox{0.96\textwidth}{!}{\begin{tikzpicture}[
  level distance=1.55cm,
  level 1/.style={sibling distance=8.0cm},
  level 2/.style={sibling distance=3.8cm},
  level 3/.style={sibling distance=1.75cm},
  edge from parent/.style={draw=gray!75,-{Latex[length=2mm]}},
  maxnode/.style={circle,draw=noteBlue,fill=noteBlue!10,minimum size=9mm,align=center,font=\scriptsize},
  minnode/.style={circle,draw=ntuRed,fill=ntuRed!9,minimum size=9mm,align=center,font=\scriptsize},
  leaf/.style={rounded corners,draw=gray!70,fill=gray!8,minimum width=7mm,minimum height=6mm,font=\scriptsize}
]
\node[maxnode] {MAX\\$A:4$}
  child { node[minnode] {MIN\\$B:4$}
    child { node[maxnode] {MAX\\$D:4$}
      child { node[leaf] {$-1$} }
      child { node[leaf] {$4$} }
    }
    child { node[maxnode] {MAX\\$E:6$}
      child { node[leaf] {$2$} }
      child { node[leaf] {$6$} }
    }
  }
  child { node[minnode] {MIN\\$C:-3$}
    child { node[maxnode] {MAX\\$F:-3$}
      child { node[leaf] {$-3$} }
      child { node[leaf] {$-5$} }
    }
    child { node[maxnode] {MAX\\$G:7$}
      child { node[leaf] {$0$} }
      child { node[leaf] {$7$} }
    }
  };
\end{tikzpicture}
}
  \caption{讲义 game tree 的 Minimax backup（价值回传）：MAX 层取最大值，MIN 层取最小值，root 最终取值 $4$。}
  \label{fig:sc3000-l03-minimax-tree}
\end{figure}

\textbf{对应例题：}\relatedexercise{SC3000-L03-E01}{Minimax Tree 的逐层回传}。

\lecturetopic{SC3000-L03-T02}{Cutoff、Evaluation、Quiescent Search 与 MCTS}

大型 game tree 无法完整展开，因此实际算法在 cutoff depth（截断深度）停止，并用 evaluation function（评估函数）近似 non-terminal state 的真实 utility：
\[
  V(s)\approx \operatorname{Eval}(s).
\]
Evaluation function 应计算高效、与 terminal utility 一致，并能合理排序局面。Othello 的简单线性评估可对双方占据的 corner、edge 与 middle 数量之差分别赋予 $3,2,1$ 的权重；它是 heuristic estimate（启发式估计），不是必然胜率。

固定深度可能产生 horizon effect（水平线效应）：剧烈后果恰好落在 cutoff 之外。Quiescent search（静态搜索）只在局面相对 quiet 时评估；若仍有明显战术变化，则继续展开到较稳定的位置。

Monte Carlo Tree Search（蒙特卡洛树搜索，MCTS）通过 selection、expansion、simulation/rollout 与 backpropagation 逐步集中计算量，并在 exploitation（利用）与 exploration（探索）之间平衡。AlphaGo 类方法以 policy network 提供 move priors，以 value network 估计 leaf value。标准缩写是 MCTS；课件中的 $f=g+h$ 只能视作宽松类比，并非标准 MCTS 更新公式。

\lecturetopic{SC3000-L03-T03}{MDP、Markov Property 与 Policy}

MDP 可写成
\[
  \boxed{\mathcal M=(S,A,P,R,\gamma)},
\]
其中 $P(s'\mid s,a)$ 是 transition model，$R(s,a,s')$ 是 reward，$\gamma$ 是 discount factor。Agent 能选择 action，但 environment 按 $P$ 随机产生 next state；因此解不是固定路径，而是 policy（策略）。Deterministic policy 写作 $\pi(s)=a$，stochastic policy 写作 $\pi(a\mid s)$。

Markov property（马尔可夫性质）为
\[
  P(S_{t+1}=s'\mid S_t=s,A_t=a,\text{history})
  =P(S_{t+1}=s'\mid S_t=s,A_t=a).
\]
它表示 current state 已包含预测未来所需的信息，不表示 transition 是 deterministic。若 state representation 丢失对未来有用的信息，Markov property 便不成立。

\lecturetopic{SC3000-L03-T04}{Return、Discount 与 Value Functions}

Discounted return（折扣回报）为
\[
  G_t=R_t+\gamma R_{t+1}+\gamma^2R_{t+2}+\cdots
      =\sum_{k=0}^{\infty}\gamma^kR_{t+k},
  \qquad 0\le\gamma<1.
\]
$\gamma$ 越小越重视 immediate reward；越接近 $1$ 越重视长期结果。Reward 有界且 $\gamma<1$ 时，无限序列的 discounted sum 仍有界。

对固定 policy $\pi$，state-value 与 action-value 分别为
\[
  V^\pi(s)=\mathbb E_\pi[G_t\mid S_t=s],
  \qquad
  Q^\pi(s,a)=\mathbb E_\pi[G_t\mid S_t=s,A_t=a],
\]
且
\[
  V^\pi(s)=\sum_{a\in A}\pi(a\mid s)Q^\pi(s,a).
\]
$V^\pi$ 表示指定 policy 的价值；只有 $V^*$ 才表示所有 policies 中的 optimal value。

\textbf{对应例题：}\relatedexercise{SC3000-L03-E02}{Game Show 的 Backward Induction}。

\lecturetopic{SC3000-L03-T05}{Bellman Optimality Equation}

Bellman equation 把 long-term value 分解为 immediate reward 与 discounted future value。固定 policy 满足
\[
  V^\pi(s)=\sum_a\pi(a\mid s)\sum_{s'}P(s'\mid s,a)
  \left[R(s,a,s')+\gamma V^\pi(s')\right].
\]
Optimal value 满足 Bellman optimality equation：
\[
  \boxed{
  V^*(s)=\max_{a\in A}\sum_{s'}P(s'\mid s,a)
  \left[R(s,a,s')+\gamma V^*(s')\right] }.
\]
其中 $\max_a$ 对应 agent 可控的 action choice；$\sum_{s'}P(s'\mid s,a)$ 对应不可控的 stochastic outcome。Optimal policy 为
\[
  \pi^*(s)=\arg\max_{a\in A}\sum_{s'}P(s'\mid s,a)
  \left[R(s,a,s')+\gamma V^*(s')\right].
\]

\begin{figure}[htbp]
  \centering
  \resizebox{0.98\textwidth}{!}{\begin{tikzpicture}[
  >=Latex,
  state/.style={rounded corners,draw=noteBlue,fill=noteBlue!9,minimum width=1.15cm,minimum height=0.7cm,font=\small},
  chance/.style={circle,draw=ntuRed,fill=ntuRed!9,minimum size=0.9cm,font=\small},
  every edge/.style={draw=gray!80,->},
  lab/.style={font=\scriptsize,fill=white,inner sep=1pt}
]
  \node[state] (s) {$s$};
  \node[chance,above right=0.7cm and 2.0cm of s] (a1) {$a_1$};
  \node[chance,below right=0.7cm and 2.0cm of s] (a2) {$a_2$};
  \node[state,above right=0.25cm and 2.2cm of a1] (s11) {$s'_1$};
  \node[state,below right=0.25cm and 2.2cm of a1] (s12) {$s'_2$};
  \node[state,above right=0.25cm and 2.2cm of a2] (s21) {$s'_3$};
  \node[state,below right=0.25cm and 2.2cm of a2] (s22) {$s'_4$};

  \draw[->,thick] (s) -- node[lab,above,sloped]{agent 选择：$\max_a$} (a1);
  \draw[->,thick] (s) -- (a2);
  \draw[->] (a1) -- node[lab,above,sloped]{$P(s'_1\mid s,a_1)$} (s11);
  \draw[->] (a1) -- node[lab,below,sloped]{$P(s'_2\mid s,a_1)$} (s12);
  \draw[->] (a2) -- node[lab,above,sloped]{$P(s'_3\mid s,a_2)$} (s21);
  \draw[->] (a2) -- node[lab,below,sloped]{$P(s'_4\mid s,a_2)$} (s22);

  \node[align=left,font=\small,right=0.9cm of s11] {
    environment 随机转移：$\sum_{s'}P(s'\mid s,a)$\\[0.25em]
    每条分支贡献：$R(s,a,s')+\gamma V(s')$
  };
\end{tikzpicture}
}
  \caption{Bellman backup 的两层结构：agent 在 action nodes 之间取最大值，environment 在 successor states 上计算期望。}
  \label{fig:sc3000-l03-mdp-backup}
\end{figure}

\textbf{对应例题：}\relatedexercise{SC3000-L03-E03}{Terminal MDP 的 Expected Action Value}。

\begin{center}
\small
\begin{tabularx}{\textwidth}{@{}lXX@{}}
  \toprule
  Problem & 下一状态的来源 & 关键 backup operation \\
  \midrule
  Ordinary search & action 通常确定 successor & 选择较优 successor \\
  Adversarial search & MAX 与 hostile MIN 交替选择 & $\max$ 与 $\min$ \\
  MDP & agent 选 action，environment 随机转移 & $\max$ 与 expectation \\
  \bottomrule
\end{tabularx}
\end{center}

\subsection{一分钟回顾}

Minimax 在 MAX/MIN 层交替回传最优与最坏结果；完整搜索不可行时，以 cutoff、evaluation、quiescent search 或 MCTS 控制计算。MDP 用 $(S,A,P,R,\gamma)$ 描述随机序贯决策，policy 指定每个 state 的 action。Value function 是 expected discounted return，Bellman optimality equation 则把一次可控的 $\max$ 与随机转移的 expectation 结合起来。Value Iteration 与 Policy Iteration 从下一讲开始。
